\documentclass[10pt,a4paper]{amsart}
\usepackage[margin=19mm]{geometry}
\usepackage{amsmath,amssymb,mathtools}
\usepackage[scr=rsfs,cal=euler]{mathalfa}
\usepackage{xfrac}
\usepackage[unicode,hidelinks,bookmarks=false]{hyperref}
\newcommand{\ie}{i.e.\ }
\newcommand{\cf}{cf.\ }
\newcommand{\resp}{respectively\ }
\newcommand{\Subsection}{\subsection}
\providecommand{\nobreakdash}{\nobreak\hbox{-}}
\setlength{\emergencystretch}{2em}
\multlinegap0pt
\usepackage{amssymb}
\usepackage{tikz}
\usepackage{xcolor}
\usepackage{mathtools}
\usepackage[normalem]{ulem}
\usepackage{dsfont}
\usepackage[shortlabels]{enumitem}
\newtheorem{lemma}{Lemma}
\newcommand{\LLn}[1][n]{{\mathsf{L}_{#1}}}
\title{Action of the Witt algebra on categorified quantum groups}
\author{Jernej Grlj, Aaron D. Lauda}
\date{Source: arXiv:2507.01877v2 (CC BY 4.0)}
\begin{document}
\maketitle

\noindent\emph{Source and licence.} Action of the Witt algebra on categorified quantum groups. Version arXiv:2507.01877v2 (CC BY 4.0), distributed by arXiv under the Creative Commons Attribution 4.0 International licence, http://creativecommons.org/licenses/by/4.0/, as recorded in the arXiv licence metadata for this exact version and shown on the abstract page https://arxiv.org/abs/2507.01877v2. Full licence notice: \texttt{licenses/arxiv-2507.01877-CC-BY-4.0.txt}.\par\smallskip

\noindent\textit{Editorial scope.} Counted unit: the article's lemma actiononfunctions (source0.tex lines 1016--1022). The article's complete original proof of it is delivered with it (source0.tex lines 1023--1058, one \texttt{proof} environment of 36 lines). No mathematics is written, repaired or re-derived: the statement and the proof are the article's own lines, in the article's own notation, and no retained line is substituted or re-flowed.  Retained uncounted as the source-local context the proof needs: lines 934--934.  Nothing was omitted from any retained span, so the package carries no omission record.  No retained line contains a citation, so the wrapper carries no bibliography and no bibliography entry.  No retained line contains a figure environment, an image-inclusion command or drawing code, so no image is delivered and the package declares no asset dependency.  Editorial formatting only, in addition: an AMS \texttt{amsart} preamble that restates the article's own declarations and macros used by the retained lines, namely the environments \texttt{lemma} on separate counters, together with the standard packages those lines need.  The retained environments print in this document as Lemma 1 in order of appearance, whereas the article prints the same items with the numbering of its own full document.  The complete native source of the article as distributed in the arXiv e-print for this exact version is bundled unchanged as \texttt{sources/180911/source0.tex}.
\subsection{Symmetric functions} \label{subsec_sympoly}

\begin{lemma}\label{actiononfunctions}
    Let $h_m$, $e_m$ and $p_m$ be the symmetric functions of degree $m$ as defined in \ref{subsec_sympoly}. Then \begin{align*}
        \LLn[n](h_m) &=  -(n+m)h_{n+m} +\sum_{j = 0}^{n-1} p_{n - j} h_{m+j} \\
        \LLn[n](e_m) &=  (-1)^{n+1} (n+m)e_{n+m} + \sum_{j = 0}^{n-1} (-1)^{j+1} p_{n - j} e_{m+j}\\
        \LLn[n](p_m) &= - m p_{n + m}  
    \end{align*}
\end{lemma}

\begin{proof}
    We will only prove the first equality; the second equality follows similarly, and the last equality is evident.
    %
    For the first equality, note that by \eqref{actiononfunctions} it suffices to show that $\LLn(h_m) = - \sum_{j = 1}^m p_{j+n} h_{m-j}$. We use the  generating function for complete symmetric functions $\sum_i h_i t^i = H(t) = \prod_{i} \frac{1}{1-x_i t}$. It follows readily that $\LLn[n]( \frac{1}{1-x_i t}) = - \frac{x_i^{n+1}t}{(1 - x_i t)^2}$, therefore
    \begin{align*}
        \LLn[n](H(t)) &= \sum_{i \geq 1} \LLn[n](h_i)t^i \\
        &= -\sum_{i \geq 1}\left(\frac{x_i^{n+1}t}{(1 - x_i t)^2} \cdot \prod_{i \neq j} (1 - x_jt)^{-1} \right) \\
        &= -\sum_{i \geq 1} \left(\frac{x_i^{n+1}t}{(1 - x_i t)} \cdot H(t) \right) \\
        &= -\sum_{i \geq 1} \sum_{k \geq 0} \left(x_i^{n+1+k}t^{k+1} \cdot H(t) \right) \\
        &= - \sum_{k \geq 1} \left(p_{k+n}t^{k} \cdot H(t) \right) \\
    \end{align*}
%     \begin{align*}
%         \LLn[n](H(t)) &= \sum_{i} \LLn[n](h_i)t^i \\
%         &= -\sum_i \left(\frac{x_i^{n+1}t}{(1 - x_i t)^2} \cdot \prod_{i \neq j} (1 - x_jt)^{-1} \right) \\
%         &= -\sum_i \left(\frac{x_i^n - x_i^n + x_i^{n+1}t}{(1 - x_i t)^2} \cdot \prod_{i \neq j} (1 - x_jt)^{-1} \right) \\
%         &= \sum_i \left(\frac{x_i^n(1 -  x_i t)}{(1 - x_i t)^2} \cdot \prod_{i \neq j} (1 - x_jt)^{-1} \right)  - \sum_i \left(\frac{x_i^n }{(1 - x_i t)^2} \cdot \prod_{i \neq j} (1 - x_jt)^{-1} \right)\\
%         &= p_n H(t) -  \sum_i \left(\frac{x_i^n }{(1 - x_i t)^2} \cdot \prod_{i \neq j} (1 - x_jt)^{-1} \right) \\
%         \end{align*}
%  where we used the fact that $p_n = \sum_i x_i^n$. We repeat the same process again, we additionally need to multiply the summation with $t \cdot \frac{1}{t}$ (and higher powers later), to get the correct numerators. 
%     \begin{align*}
%         &= p_n H(t) -   \frac{1}{t}\sum_i \left(\frac{x_i^{n-1} - x_i^{n-1} + t x_i^n }{(1 - x_i t)^2} \cdot \prod_{i \neq j} (1 - x_jt)^{-1} \right) \\
%         &= p_n H(t) + \frac{1}{t} p_{n-1}H(t) -   \frac{1}{t}\sum_i \left(\frac{x_i^{n-1}}{(1 - x_i t)^2} \cdot \prod_{i \neq j} (1 - x_jt)^{-1} \right) \\
%         &= \sum_{j = 0}^{n-1} t^{-j} p_{n - j} H(t) - \frac{1}{t^{n-1}}\sum_i \left(\frac{x_i}{(1 - x_i t)^2} \cdot \prod_{i \neq j} (1 - x_jt)^{-1} \right) \\
%         \end{align*}
%         One last trick is to realize the terms in the summation as partial derivatives in $t$.
%         \begin{align*}
%         &= \sum_{j = 0}^{n-1} t^{-j} p_{n - j} H(t) - \frac{1}{t^{n-1}} \frac{\partial}{\partial t} \prod_{i } (1 - x_jt)^{-1}  \\
%         &= \sum_{j = 0}^{n-1} t^{-j} p_{n - j} H(t) - \frac{1}{t^{n-1}} \frac{\partial}{\partial t} \sum_i h_i t^i \\
%         &= \sum_{j = 0}^{n-1} t^{-j} p_{n - j} H(t) - \sum_i i h_i t^{i-n} 
%     \end{align*}
% So we observe.
% \begin{align}
%  \LLn[n](h_m) &= -(n+m)h_{n+m} +\sum_{j = 0}^{n-1} p_{n - j} h_{m+j}
% \end{align}
This completes the proof after we expand $H(t)$.
\end{proof}

\end{document}
